% Research-programme notes, read by notesmap and buildable as a PDF.
% One programme per file under programmes/ keeps edits and locks small:
% two people can work on different programmes without touching the same file.
\documentclass[11pt, a4paper]{article}
\usepackage[margin=2.5cm]{geometry}
\usepackage[hidelinks]{hyperref}
\usepackage{notesmap}

\begin{document}
\tableofcontents
\clearpage

\begin{programnote}{Example Lab -- what the group builds}{Example}
\location{51.50}{-0.12}{London}
% \thumb{img/example.png}   % put the image in img/ and uncomment

\paragraph{What it is}
One or two sentences on who the group is and what they build.

\paragraph{Why it matters here}
Why this programme matters to your work. A cross-reference to another
programme's label (prog:its-key) or paper's label (note:its-key), made with
the usual ref command, becomes an arrow on the map, captioned by its sentence.

\begin{papernote}{Author et al.\ 2024 -- Short title}{Author2024}
Journal 1:2, \doi{10.0000/example}.

\paragraph{Summary}
What the paper did, in plain terms. \gap{What is still to be checked.}

\end{papernote}

\end{programnote}

\end{document}
